% !TEX root = ../main.tex

\chapter{实验算法伪代码}\label{ch:pseudocode}
  本附录给出本文实验中实际采用的求解伪代码.
  为了便于与正文中的方法分析对应, 这里不再强调具体实现细节,
  而是统一使用“路径采样, 特征构造, 参数求解, 测试评估”的算法表述.

  \section{线性方程求解流程}\label{sec:app_linear_method}
    \begin{algorithm}[htbp]
      \SetAlgoLined
      \KwData{训练路径, 布朗增量, 随机特征, 线性方程离散系数, 岭正则化参数 $\lambda_{\mathrm{r}}$}
      \KwResult{$y_{0}$, $\alpha$}
      对每条路径计算终端表达式中的系数 $\beta_{i}$, $\Gamma_{i}$, $\eta_{i}$\;
      将所有路径拼接为设计矩阵 $A$ 与观测向量 $B$\;
      令参数向量
      \begin{equation*}
        \vartheta =
        \begin{bmatrix}
          y_{0} \\
          \operatorname{vec}(\alpha)
        \end{bmatrix},
      \end{equation*}
      并记 $p = 1 + dH$\;
      构造
      \begin{equation*}
        R =
        \begin{bmatrix}
          0 & 0 \\
          0 & I_{dH}
        \end{bmatrix}
      \end{equation*}
      用于仅对 $\alpha$ 正则化\;
      \eIf{$S\ge p$}{
        求解原始岭正则化系统
        $
          \left(A^{\top}A+\lambda_{\mathrm{r}}R^{\top}R\right)\vartheta
          =
          A^{\top}B
        $
        得到参数 $\theta$\;
      }{
        求解对偶岭正则化系统
        $
          \vartheta = A^{\top}
          \left(AA^{\top} + \lambda_{\mathrm{r}} I_S\right)^{-1}B\;
        $
        得到参数 $\theta$\;
      }
      解包 $\theta$ 得到参数 $y_{0}$ 与 $\alpha$\;
      返回 $y_{0}$ 与 $\alpha$\;
      \caption{线性方程的求解流程}
      \label{alg:linear}
    \end{algorithm}

    \newpage

  \section{半线性方程求解流程}\label{sec:app_semilinear_method}
    \begin{algorithm}[htbp]
      \SetAlgoLined
      \KwData{训练路径, 布朗增量, 随机特征, 初始参数 $\theta^{(0)}$, 初始阻尼参数 $\lambda^{(0)}$}
      \KwResult{$y_{0}$, $\alpha$}
      令 $m \leftarrow 0$\;
      \While{未满足停止准则}{
        根据当前参数 $\theta^{(m)}$ 构造所有 $Z_{k}^{(i)}$\;
        前向推进离散 BSDE, 计算终端残差向量 $r(\theta^{(m)})$\;
        计算雅可比矩阵 $J(\theta^{(m)})$\;
        求解阻尼子问题
        \begin{equation*}
          \bigl(J^{\top} J + \lambda^{(m)} I\bigr)\delta^{(m)}
          =
          -J^{\top} r;
        \end{equation*}
        令试探参数 $\theta_{\mathrm{trial}} = \theta^{(m)} + \delta^{(m)}$\;
        计算试探参数对应的目标函数值\;
        \eIf{目标函数下降}{
          接受该步, 令 $\theta^{(m + 1)} = \theta_{\mathrm{trial}}$\;
          适当减小阻尼参数 $\lambda^{(m)}$\;
          $m \leftarrow m + 1$\;
        }{
          拒绝该步并增大阻尼参数, 重新计算当前迭代的试探步\;
        }
      }
      从最终参数向量中拆分出 $y_{0}$ 与 $\alpha$\;
      返回 $y_{0}$ 与 $\alpha$\;
      \caption{半线性方程的 LM 求解流程}
      \label{alg:semilinear}
    \end{algorithm}

    \newpage

  \section{测试误差评估流程}\label{sec:app_test_method}
    \begin{algorithm}[htbp]
      \SetAlgoLined
      \KwData{训练得到的 $y_{0}$, $\alpha$}
      \KwResult{独立测试样本上的终端残差 RMSE}
      重新生成独立测试路径 $\{X_{k}^{(i)}, \Delta W_{k}^{(i)}\}$\;
      在测试路径上构造随机特征 $\{H_{k}^{(i)}\}$\;
      用固定参数 $\alpha$ 构造 $Z_{k}^{(i)} = \alpha H_{k}^{(i)}$\;
      沿测试路径前向推进离散 BSDE, 得到所有终端值 $Y_{N}^{(i)}$\;
      计算
      \begin{equation*}
        \operatorname{RMSE}_{\text{test}}
        =
        \sqrt{
          \frac{1}{S}
          \sum_{i = 1}^{S}
          \bigl(
          Y_{N}^{(i)} - g(X_{N}^{(i)})
          \bigr)^{2}
        };
      \end{equation*}
      返回测试 RMSE\;
      \caption{测试误差评估流程}
      \label{alg:test}
    \end{algorithm}
